%% methodology.tex
%%

%% ==============================
\section{Methodology}
\label{sec:Methodology}
%% ==============================

This section describes the experimental methodology. Formal CA2 specifies GCT/Alibaba LSTM training, a Kubernetes-API adaptive scaler, and evaluation on AWS-hosted Kubernetes with the metric suite in Table~\ref{tab:formal-metrics-intent}. \textbf{Live AWS Kubernetes is not run this pass.} The runnable formal path uses a public Google Cluster Data v1 (2010) slice (CC-BY; SHA1-verified) for LSTM windows, a NumPy LSTM, and a dry-run \texttt{apps/v1} Scale client versus an \texttt{autoscaling/v2} HPA object. GCT 2011/2019 and Alibaba requests fail closed (\texttt{DATA\_GAPS.md}). A synthetic sine-spike generator remains as an explicit \textbf{proxy} only.

\begin{table}[htp]
\centering
\caption{Formal CA2 metrics and evidence in this artefact}
\label{tab:formal-metrics-intent}
\begin{tabular}{lp{8cm}}
\hline
\textbf{Metric} & \textbf{This pass} \\
\hline
Prediction MAE, RMSE & TRACE on GCT v1 job-held-out windows; cluster series is short ($\sim$7\,h / 5\,min) \\
CPU / memory utilisation & SIMULATED from load / (replicas $\times$ assumed capacity) \\
Response time, throughput & SIMULATED queue / cap model; not live RPS traces \\
Scaling latency & SIMULATED as one control interval (300\,s bins); not kubelet Ready time \\
Infrastructure cost & SIMULATED (assumed \$0.04/pod-hour); not AWS billing \\
SLA / availability & SIMULATED from capacity mapping \\
\hline
\end{tabular}
\end{table}

\subsection{Trace-derived workload (formal path)}

The committed public slice is Google Cluster Data version 1~\cite{GCT2010} (7-hour Borg cell). Rows are 5-minute task reports of normalised CPU and memory. Aggregation (no imputed gaps): cluster CPU is $\sum$ \texttt{NrmlTaskCores} per \texttt{Time}; LSTM windows are per-job series with at least 12 bins (\texttt{data/traces/PROVENANCE.md}). This is cluster-trace data. It is not the 2011 29-day cell, the 2019 eight-cell Borg traces, or Alibaba.

\subsection{Synthetic Workload Generation (PROXY only)}


A synthetic cyclical workload (\texttt{src/data/workload\_simulator.py}) of 500 timesteps is generated for the \textbf{proxy} NimbusGuard-framed experiment only. It is not a silent fallback for \texttt{--dataset gct*} or \texttt{alibaba}. The workload combines three components:

\paragraph{Base Load} A sinusoidal pattern with amplitude 20 and period 100 timesteps, representing diurnal or weekly usage cycles common in cloud applications:
\[
\text{base}(t) = 30 + 20 \sin\left(\frac{2\pi t}{100}\right)
\]

\paragraph{Periodic Spikes} Traffic spikes of magnitude 15-25 (uniformly random) occur every 75-85 timesteps (uniformly random), lasting 5 timesteps each. These spikes simulate flash crowds, batch jobs, or marketing campaigns that arrive on top of the baseline pattern.

\paragraph{Gaussian Noise} Zero-mean Gaussian noise with standard deviation 3.0 is added at each timestep to represent the inherent stochasticity of real workloads -- the kind of high-frequency variation that aggressive proactive scaling can overreact to.

The final workload at timestep $t$ is:
\[
W(t) = \max\left(5, \text{base}(t) + \text{spike}(t) + \mathcal{N}(0, 3^2)\right)
\]
The floor of 5 ensures that workload never drops below a minimal baseline.

\subsection{Predictor Architecture and Training}

\paragraph{Formal path} A univariate LSTM (lookback 12, hidden 16) is trained on GCT v1 job-level CPU windows with a held-out-job split (no window leakage). Runtime is NumPy BPTT; TensorFlow is fail-closed if requested. Cluster-level walk-forward MAE/RMSE is also reported but $n$ is small (76 five-minute bins).

\paragraph{Proxy path} Proactive proxy policies share a scikit-learn \texttt{MLPRegressor} (64/32). This is not the binding LSTM.

\paragraph{Proxy training procedure} For each random seed $s \in \{42, 43, 44, 45, 46\}$:
\begin{enumerate}
  \item Generate a fresh 500-step workload $W_s(t)$ using seed $s$.
  \item Construct a supervised training dataset from the first 300 timesteps: for each timestep $t \in [10, 300]$, the input is the 10-step lookback window $[W_s(t-10), \ldots, W_s(t-1)]$ and the target is $W_s(t)$.
  \item Train the predictor on this dataset for 1000 epochs with early stopping (patience 50) on a held-out 20\% validation split.
  \item Evaluate all three policies (Reactive HPA, Aggressive PAKS, Stability-Aware PAKS) on the \emph{full} 500-step workload $W_s(t)$, using the trained predictor where applicable.
\end{enumerate}

This setup ensures that:
\begin{itemize}
  \item All three policies are compared on identical workloads per seed, isolating the effect of the scaling policy.
  \item The predictor is trained only on the first 60\% of the sequence, testing its ability to generalize to unseen patterns (spikes in the latter 40\%).
  \item Each seed gets a fresh workload and fresh predictor, so results average over different workload realizations.
\end{itemize}

\subsection{Scaling Policies}

\paragraph{Reactive HPA (binding baseline)} Matches the Kubernetes Horizontal Pod Autoscaler v2 CPU-utilization replica formula on \emph{observed} load, emitting \texttt{autoscaling/v2} status objects and \texttt{PATCH} \texttt{apps/v1} \texttt{/scale?dryRun=All}. One simulated control-interval delay. Not a live metrics-server scrape.

\paragraph{PAKS adaptive engine (formal stub)} Uses the LSTM next-step forecast, converts predicted load to a replica count, and emits the same Scale subresource shape (dry-run). Replica application in the loop is simulated (no kubelet).

The following three policies are the \textbf{proxy} simulator (MLP, synthetic load) and are not the CA2 comparison contract.

\paragraph{Reactive HPA (proxy discrete-time mimic)} Mimics HPA with a one-step delay on synthetic load (\texttt{capacity\_per\_pod} = 100 in code; older draft text used 10 --- the CSV is bound to the code).

\paragraph{Aggressive PAKS (Baseline 2)} At each timestep $t$, use the predictor to forecast $\hat{W}(t+1)$. Apply a 5\% safety buffer and convert directly to a target pod count with no smoothing or cooldown:
\[
\text{pods}_{\text{agg}}(t) = \left\lceil \frac{1.05 \cdot \hat{W}(t+1)}{\text{capacity\_per\_pod}} \right\rceil
\]
This reproduces the core behavior of NimbusGuard's DQN+LSTM agent: proactive, but reacting to every raw forecast including noise. This is the "baseline reproduction" that demonstrates the agility-instability trade-off.

\paragraph{Stability-Aware PAKS (Proposed)} Built on the exact same predictor as Aggressive PAKS, but wrapped in a \texttt{StabilityAwareController} that applies two stability mechanisms before executing a scaling action:

\begin{enumerate}
  \item \textbf{Exponential Moving Average (EMA)} of predictions: instead of acting on the raw forecast $\hat{W}(t+1)$, maintain a smoothed estimate:
  \[
  \text{smoothed}(t) = \alpha \cdot \hat{W}(t+1) + (1 - \alpha) \cdot \text{smoothed}(t-1)
  \]
  where $\alpha = 0.4$. This filters high-frequency noise while retaining responsiveness to genuine trends.
  
  \item \textbf{Hysteresis and Cooldown}: Compute the desired pod count from the smoothed prediction:
  \[
  \text{desired}(t) = \left\lceil \frac{1.05 \cdot \text{smoothed}(t)}{\text{capacity\_per\_pod}} \right\rceil
  \]
  A scaling action only fires if \emph{both} of the following hold:
  \begin{itemize}
    \item $|\text{desired}(t) - \text{current\_pods}(t)| > 1$ (hysteresis: at least 2 pods difference)
    \item $t - t_{\text{last\_scale}} \geq 2$ (cooldown: at least 2 timesteps since last action)
  \end{itemize}
  Otherwise, the pod count remains unchanged.
\end{enumerate}

This is a direct, minimal implementation of the ``calmer design'' NimbusGuard's future-work section names but does not build -- exponential smoothing and hysteresis/cooldown are the classical stability mechanisms documented in the literature (see Section~\ref{sec:rw}), applied here as simple rules rather than a second learned agent.

\begin{figure}[htbp]
\centering
\begin{tikzpicture}[
  block/.style={rectangle, draw, fill=blue!10, text width=2.5cm, text centered, rounded corners, minimum height=0.8cm, font=\small},
  process/.style={rectangle, draw, fill=green!10, text width=2.5cm, text centered, rounded corners, minimum height=0.8cm, font=\small},
  decision/.style={diamond, draw, fill=yellow!10, text width=1.8cm, text centered, minimum height=0.8cm, font=\scriptsize},
  line/.style={draw, -Latex},
  node distance=1.2cm and 1.5cm
]

% Reactive HPA (Left column)
\node[block] (hpa-workload) {Observed\\Workload $W(t-1)$};
\node[process, below=of hpa-workload] (hpa-calc) {Direct\\Calculation};
\node[block, below=of hpa-calc] (hpa-pods) {Target\\Pods};

\draw[line] (hpa-workload) -- (hpa-calc);
\draw[line] (hpa-calc) -- (hpa-pods);

\node[above=0.3cm of hpa-workload, font=\bfseries] {Reactive HPA};

% Aggressive PAKS (Middle column)
\node[block, right=2.5cm of hpa-workload] (agg-workload) {Historical\\Workload};
\node[process, below=of agg-workload] (agg-predictor) {ML Predictor\\(MLP)};
\node[block, below=of agg-predictor] (agg-forecast) {Raw Forecast\\$\hat{W}(t+1)$};
\node[process, below=of agg-forecast] (agg-buffer) {+5\% Safety\\Buffer};
\node[block, below=of agg-buffer] (agg-pods) {Target\\Pods};

\draw[line] (agg-workload) -- (agg-predictor);
\draw[line] (agg-predictor) -- (agg-forecast);
\draw[line] (agg-forecast) -- (agg-buffer);
\draw[line] (agg-buffer) -- (agg-pods);

\node[above=0.3cm of agg-workload, font=\bfseries] {Aggressive PAKS};

% Stability-Aware PAKS (Right column)
\node[block, right=2.5cm of agg-workload] (sa-workload) {Historical\\Workload};
\node[process, below=of sa-workload] (sa-predictor) {ML Predictor\\(MLP)};
\node[block, below=of sa-predictor] (sa-forecast) {Raw Forecast\\$\hat{W}(t+1)$};
\node[process, below=of sa-forecast] (sa-ema) {EMA Smoothing\\($\alpha=0.4$)};
\node[process, below=of sa-ema] (sa-buffer) {+5\% Safety\\Buffer};
\node[decision, below=of sa-buffer] (sa-hyst) {$\Delta > 1$ pod\\AND\\cooldown?};
\node[block, below=0.8cm of sa-hyst] (sa-pods) {Target\\Pods};

\draw[line] (sa-workload) -- (sa-predictor);
\draw[line] (sa-predictor) -- (sa-forecast);
\draw[line] (sa-forecast) -- (sa-ema);
\draw[line] (sa-ema) -- (sa-buffer);
\draw[line] (sa-buffer) -- (sa-hyst);
\draw[line] (sa-hyst) -- node[right, font=\scriptsize] {yes} (sa-pods);
\draw[line] (sa-hyst.west) -- ++(-0.8,0) |- node[left, pos=0.25, font=\scriptsize] {no: hold} (sa-pods.west);

\node[above=0.3cm of sa-workload, font=\bfseries] {Stability-Aware PAKS};

% Legend
\node[below=1.2cm of hpa-pods, font=\scriptsize\itshape, text width=10cm, text centered] {
  \textbf{Key difference:} Reactive HPA reacts to past workload; both PAKS policies predict future demand. 
  Stability-Aware adds smoothing + hysteresis/cooldown to prevent overreaction to noise.
};

\end{tikzpicture}
\caption{Controller architecture dataflow for the three scaling policies. Reactive HPA uses only observed workload with a delay. Aggressive PAKS applies raw predictions directly. Stability-Aware PAKS wraps the same predictor in exponential smoothing and hysteresis/cooldown gates to filter transient fluctuations.}
\label{fig:controller_architecture}
\end{figure}


\subsection{Evaluation Metrics}

Four metrics capture different facets of autoscaling performance:

\paragraph{SLA Violations} Count of timesteps where allocated pods fall below the absolute minimum required to serve the actual workload without throttling:
\[
\text{violations} = \sum_{t=1}^{500} \mathbb{1}\left[\text{pods}(t) < \left\lceil \frac{W(t)}{\text{capacity\_per\_pod}} \right\rceil\right]
\]
Lower is better. This is the primary performance metric that motivates proactive scaling.

\paragraph{Over-provisioning \%} Total allocated pod-steps in excess of the minimum required, as a percentage:
\[
\text{over\_prov} = 100 \cdot \frac{\sum_{t=1}^{500} \max\left(0, \text{pods}(t) - \left\lceil \frac{W(t)}{\text{capacity\_per\_pod}} \right\rceil\right)}{\sum_{t=1}^{500} \left\lceil \frac{W(t)}{\text{capacity\_per\_pod}} \right\rceil}
\]
Lower is better. This quantifies resource waste and cost.

\paragraph{Scaling Events} Count of timesteps where the pod count changes from the previous timestep:
\[
\text{events} = \sum_{t=2}^{500} \mathbb{1}[\text{pods}(t) \neq \text{pods}(t-1)]
\]
Lower is better. This is directly comparable to NimbusGuard's own ``Total Scaling Events'' metric and quantifies system instability.

\paragraph{Pod Count Volatility} Standard deviation of step-to-step pod-count changes:
\[
\text{volatility} = \sqrt{\frac{1}{499} \sum_{t=2}^{500} \left(\text{pods}(t) - \text{pods}(t-1)\right)^2}
\]
Lower is better. This is a finer-grained stability measure that captures both the frequency and magnitude of scaling actions, penalizing large swings more than small adjustments.

All code and results are in \texttt{paks-framework/}. Running \texttt{python scripts/train\_and\_evaluate.py} reproduces all experimental results from scratch.

\subsection{Statistical Rigor}

Results are reported as mean $\pm$ standard deviation across 5 independent seeds (42--46). Each seed generates a different workload realization and trains a fresh predictor, ensuring that findings are not artifacts of a single workload pattern. Standard deviations are reported for all metrics to quantify variability. Differences between policies are considered meaningful when they exceed one standard deviation of the more variable policy.
